\chapter{Geometría}
\label{ch:geometria}

La geometría analítica permite traducir problemas geométricos a ecuaciones y, recíprocamente, interpretar geométricamente sistemas de ecuaciones. En ingeniería esta doble lectura es especialmente útil: la posición de un eje, la orientación de una superficie, la trayectoria ideal de una herramienta, la dirección de una fuerza o la separación mínima entre dos elementos pueden describirse mediante puntos, vectores, rectas y planos.

En el capítulo anterior hemos trabajado con espacios vectoriales, producto escalar y ortogonalidad. Ahora utilizaremos esas herramientas en $\R^2$ y $\R^3$ para estudiar la \emph{geometría afín} ---incidencia, paralelismo e intersecciones--- y la \emph{geometría euclídea} ---distancias, ángulos y perpendicularidad---. La intención es pasar con soltura entre diferentes registros: una descripción verbal de una situación, su representación vectorial y las ecuaciones que permiten resolverla.

\section{Puntos y vectores en el plano y el espacio}
\label{sec:geometria-vectores}

\subsection{Coordenadas, vectores y distancia entre puntos}

Fijado un sistema de referencia cartesiano, un punto del plano se representa mediante dos coordenadas,
\[
P=(p_1,p_2)\in\R^2,
\]
y un punto del espacio mediante tres,
\[
P=(p_1,p_2,p_3)\in\R^3.
\]
Conviene distinguir conceptualmente \emph{puntos} y \emph{vectores}. Un punto representa una posición; un vector representa un desplazamiento, una dirección o, en muchas aplicaciones, una magnitud física con módulo y dirección.

Dados dos puntos $P$ y $Q$, el vector que lleva de $P$ a $Q$ es
\[
\overrightarrow{PQ}=Q-P.
\]
En $\R^3$, si
\[
P=(p_1,p_2,p_3),\qquad Q=(q_1,q_2,q_3),
\]
entonces
\[
\overrightarrow{PQ}
=(q_1-p_1,q_2-p_2,q_3-p_3).
\]

\begin{definition}[Norma de un vector]
Para $v=(v_1,\ldots,v_n)\in\R^n$, su \emph{norma euclídea} es
\[
\|v\|=\sqrt{v_1^2+\cdots+v_n^2}.
\]
\end{definition}

La norma mide la longitud del vector. Si $v\neq0$, el vector
\[
\widehat v=\frac{v}{\|v\|}
\]
tiene norma $1$ y conserva la dirección y el sentido de $v$. Lo denominamos \emph{vector unitario} asociado a $v$.

\begin{definition}[Distancia entre dos puntos]
La distancia entre dos puntos $P,Q\in\R^n$ es
\[
d(P,Q)=\|Q-P\|.
\]
\end{definition}

Así, en el espacio,
\[
d(P,Q)
=\sqrt{(q_1-p_1)^2+(q_2-p_2)^2+(q_3-p_3)^2}.
\]

\begin{example}
El extremo de un actuador se desplaza desde
\[
P=(1,-2,0.5)\ \mathrm{m}
\]
hasta
\[
Q=(4,2,2.5)\ \mathrm{m}.
\]
El desplazamiento es
\[
\overrightarrow{PQ}=(3,4,2)\ \mathrm{m},
\]
y la distancia recorrida en línea recta entre ambas posiciones es
\[
\|\overrightarrow{PQ}\|
=\sqrt{3^2+4^2+2^2}
=\sqrt{29}\ \mathrm{m}.
\]
La dirección unitaria del desplazamiento es
\[
\widehat{PQ}=\frac{1}{\sqrt{29}}(3,4,2).
\]
Si el movimiento se realizara exactamente siguiendo esa dirección, cualquier posición de la trayectoria podría escribirse como
\[
P+t\widehat{PQ},\qquad t\geq0.
\]
\end{example}

La expresión anterior anticipa la ecuación vectorial de una recta: una posición inicial más un múltiplo de una dirección.

\subsection{Producto escalar, ángulos y proyecciones}

Recordemos que, dados $u,v\in\R^n$, su producto escalar es
\[
u\cdot v=u_1v_1+\cdots+u_nv_n.
\]
Además,
\[
u\cdot v=\|u\|\,\|v\|\cos\theta,
\]
donde $\theta\in[0,\pi]$ es el ángulo que forman los vectores no nulos $u$ y $v$. Por tanto,
\[
\cos\theta=\frac{u\cdot v}{\|u\|\,\|v\|}.
\]

\begin{proposition}[Criterio de ortogonalidad]
Dos vectores $u$ y $v$ son ortogonales si y sólo si
\[
u\cdot v=0.
\]
\end{proposition}

Esta igualdad permite traducir perpendicularidad geométrica en una ecuación algebraica.

\begin{definition}[Proyección ortogonal sobre una dirección]
Sean $u,v\in\R^n$, con $v\neq0$. La proyección ortogonal de $u$ sobre la dirección de $v$ es
\[
\operatorname{proj}_v(u)
=\frac{u\cdot v}{v\cdot v}\,v.
\]
\end{definition}

El vector $u$ puede descomponerse como
\[
u=u_{\parallel}+u_{\perp},
\]
donde
\[
u_{\parallel}=\operatorname{proj}_v(u),
\qquad
u_{\perp}=u-\operatorname{proj}_v(u),
\]
y $u_{\perp}\cdot v=0$.

\begin{solvedexercise}
Una fuerza
\[
F=(120,40,-20)\ \mathrm{N}
\]
actúa sobre un mecanismo cuyo desplazamiento permitido tiene la dirección
\[
d=(2,1,2).
\]
Calcular la componente vectorial de $F$ que actúa en la dirección del movimiento y la componente perpendicular.
\tcblower
\textbf{Resolución.}\par\smallskip
La proyección de $F$ sobre $d$ es
\[
F_{\parallel}
=\frac{F\cdot d}{d\cdot d}\,d.
\]
Como
\[
F\cdot d=120\cdot2+40\cdot1-20\cdot2=240,
\]
y
\[
d\cdot d=2^2+1^2+2^2=9,
\]
obtenemos
\[
F_{\parallel}=\frac{240}{9}(2,1,2)
=\left(\frac{160}{3},\frac{80}{3},\frac{160}{3}\right)\ \mathrm{N}.
\]
Por tanto,
\[
F_{\perp}=F-F_{\parallel}
=\left(\frac{200}{3},\frac{40}{3},-\frac{220}{3}\right)\ \mathrm{N}.
\]
Puede comprobarse que
\[
F_{\perp}\cdot d=0.
\]
La descomposición separa la parte de la fuerza que contribuye al movimiento permitido de la parte que actúa transversalmente sobre el mecanismo.
\end{solvedexercise}

En $\R^3$ también es útil describir una dirección mediante los ángulos que forma con los ejes coordenados. Si $v=(v_1,v_2,v_3)\neq0$ y $\alpha,\beta,\gamma$ son los ángulos con los ejes $OX$, $OY$ y $OZ$, respectivamente, entonces
\[
\cos\alpha=\frac{v_1}{\|v\|},
\qquad
\cos\beta=\frac{v_2}{\|v\|},
\qquad
\cos\gamma=\frac{v_3}{\|v\|},
\]
y, en consecuencia,
\[
\cos^2\alpha+\cos^2\beta+\cos^2\gamma=1.
\]

\subsection{Producto vectorial y producto mixto}

En $\R^3$ disponemos de una operación adicional que produce un vector perpendicular a otros dos.

\begin{definition}[Producto vectorial]
Sean
\[
u=(u_1,u_2,u_3),\qquad v=(v_1,v_2,v_3).
\]
Su \emph{producto vectorial} es
\[
u\times v=
\begin{vmatrix}
\mathbf i&\mathbf j&\mathbf k\\
u_1&u_2&u_3\\
v_1&v_2&v_3
\end{vmatrix},
\]
es decir,
\[
u\times v
=
(u_2v_3-u_3v_2,\,
 u_3v_1-u_1v_3,\,
 u_1v_2-u_2v_1).
\]
\end{definition}

El producto vectorial satisface:
\begin{itemize}
\item $u\times v$ es ortogonal a $u$ y a $v$;
\item $v\times u=-(u\times v)$;
\item $u\times v=0$ si y sólo si $u$ y $v$ son linealmente dependientes;
\item
\[
\|u\times v\|=\|u\|\,\|v\|\sin\theta.
\]
\end{itemize}

La última propiedad tiene una interpretación geométrica inmediata:
\[
\text{Área del paralelogramo generado por }u,v
=\|u\times v\|,
\]
y, por tanto,
\[
\text{Área del triángulo}
=\frac12\|u\times v\|.
\]

\begin{example}
Tres puntos de una placa son
\[
A=(1,0,0),\qquad B=(3,1,0),\qquad C=(2,0,2).
\]
Los vectores
\[
\overrightarrow{AB}=(2,1,0),
\qquad
\overrightarrow{AC}=(1,0,2)
\]
determinan la placa. Un vector perpendicular a ella es
\[
\overrightarrow{AB}\times\overrightarrow{AC}
=
\begin{vmatrix}
\mathbf i&\mathbf j&\mathbf k\\
2&1&0\\
1&0&2
\end{vmatrix}
=(2,-4,-1).
\]
Este vector normal será precisamente el que utilicemos más adelante para obtener la ecuación cartesiana del plano que contiene la placa.
\end{example}

\begin{definition}[Producto mixto]
Dados $u,v,w\in\R^3$, se denomina \emph{producto mixto} al número
\[
[u,v,w]=u\cdot(v\times w).
\]
\end{definition}

Se verifica
\[
[u,v,w]
=
\det
\begin{pmatrix}
u_1&u_2&u_3\\
v_1&v_2&v_3\\
w_1&w_2&w_3
\end{pmatrix}.
\]
Su valor absoluto es el volumen del paralelepípedo generado por $u,v,w$:
\[
V=|u\cdot(v\times w)|.
\]
Además,
\[
u,v,w\text{ son coplanarios}
\quad\Longleftrightarrow\quad
u\cdot(v\times w)=0.
\]

\begin{solvedexercise}
Una fuerza $F=(0,150,60)$ N actúa en el punto cuyo vector de posición respecto de un eje de referencia es
\[
r=(0.4,0,0.2)\ \mathrm{m}.
\]
Calcular el momento de la fuerza respecto del origen,
\[
M=r\times F.
\]
\tcblower
\textbf{Resolución.}\par\smallskip
Calculamos
\[
M=
\begin{vmatrix}
\mathbf i&\mathbf j&\mathbf k\\
0.4&0&0.2\\
0&150&60
\end{vmatrix}
=(-30,-24,60)\ \mathrm{N\,m}.
\]
El vector momento es perpendicular al plano generado por $r$ y $F$. Su módulo es
\[
\|M\|=\sqrt{(-30)^2+(-24)^2+60^2}
=\sqrt{5076}\ \mathrm{N\,m}.
\]
\end{solvedexercise}

\section{Rectas y planos}
\label{sec:rectas-planos}

\subsection{Rectas en el plano y en el espacio}

Una recta queda determinada por un punto y una dirección.

\begin{definition}[Ecuación vectorial de una recta]
Sea $A\in\R^n$ y sea $v\in\R^n$, $v\neq0$. La recta que pasa por $A$ y tiene vector director $v$ es
\[
r:\quad X=A+tv,\qquad t\in\R.
\]
\end{definition}

En el plano, si
\[
A=(a_1,a_2),\qquad v=(v_1,v_2),
\]
las ecuaciones paramétricas son
\[
\begin{cases}
x=a_1+tv_1,\\
y=a_2+tv_2,
\end{cases}
\qquad t\in\R.
\]
En el espacio,
\[
\begin{cases}
x=a_1+tv_1,\\
y=a_2+tv_2,\\
z=a_3+tv_3.
\end{cases}
\]

Cuando las componentes del vector director son no nulas, puede eliminarse el parámetro y escribirse la forma continua o simétrica
\[
\frac{x-a_1}{v_1}
=
\frac{y-a_2}{v_2}
\]
en el plano, y
\[
\frac{x-a_1}{v_1}
=
\frac{y-a_2}{v_2}
=
\frac{z-a_3}{v_3}
\]
en el espacio. Si alguna componente del vector director es cero, la coordenada correspondiente permanece constante y no debe dividirse por cero.

En $\R^2$ una recta también puede expresarse mediante una ecuación cartesiana
\[
r:\quad Ax+By+C=0,
\qquad (A,B)\neq(0,0).
\]
El vector
\[
n=(A,B)
\]
es normal a la recta. Por tanto, un vector director puede elegirse como
\[
v=(-B,A).
\]
Si $B\neq0$, la ecuación puede escribirse
\[
y=-\frac ABx-\frac CB.
\]
Así, la pendiente es $m=-A/B$.

\begin{example}
Hallar una ecuación de la recta del plano que pasa por
\[
P=(1,-2),\qquad Q=(5,1).
\]
Un vector director es
\[
v=Q-P=(4,3).
\]
Por tanto,
\[
(x,y)=(1,-2)+t(4,3),
\]
o, en forma paramétrica,
\[
\begin{cases}
x=1+4t,\\
y=-2+3t.
\end{cases}
\]
Un vector normal es $n=(3,-4)$, ya que
\[
(3,-4)\cdot(4,3)=0.
\]
La ecuación cartesiana es entonces
\[
3(x-1)-4(y+2)=0,
\]
es decir,
\[
\boxed{3x-4y-11=0.}
\]
\end{example}

Una recta del espacio también puede describirse como intersección de dos planos independientes:
\[
r:\quad
\begin{cases}
A_1x+B_1y+C_1z+D_1=0,\\
A_2x+B_2y+C_2z+D_2=0.
\end{cases}
\]
Si $n_1=(A_1,B_1,C_1)$ y $n_2=(A_2,B_2,C_2)$ no son paralelos, un vector director de la recta es
\[
v=n_1\times n_2.
\]

\begin{solvedexercise}
El eje de una tubería recta pasa por los puntos
\[
A=(1,2,0),\qquad B=(4,-1,6).
\]
Obtener sus ecuaciones vectorial, paramétricas y, cuando sea posible, continuas.
\tcblower
\textbf{Resolución.}\par\smallskip
Un vector director es
\[
v=B-A=(3,-3,6)=3(1,-1,2).
\]
Podemos utilizar la dirección simplificada $(1,-1,2)$. La ecuación vectorial es
\[
\boxed{(x,y,z)=(1,2,0)+t(1,-1,2).}
\]
Las ecuaciones paramétricas son
\[
\boxed{
\begin{cases}
x=1+t,\\
y=2-t,\\
z=2t.
\end{cases}}
\]
Como ninguna componente del vector director es nula,
\[
\boxed{x-1=-(y-2)=\frac z2.}
\]
Las tres expresiones describen exactamente el mismo conjunto de puntos.
\end{solvedexercise}

\subsection{Planos en el espacio}

Un plano puede determinarse mediante un punto y dos direcciones linealmente independientes.

\begin{definition}[Ecuación vectorial de un plano]
Sean $A\in\R^3$ y $u,v\in\R^3$ linealmente independientes. El plano que pasa por $A$ y tiene direcciones $u$ y $v$ es
\[
\pi:\quad X=A+su+tv,
\qquad s,t\in\R.
\]
\end{definition}

Si
\[
A=(a_1,a_2,a_3),\quad
u=(u_1,u_2,u_3),\quad
v=(v_1,v_2,v_3),
\]
obtenemos las ecuaciones paramétricas
\[
\begin{cases}
x=a_1+su_1+tv_1,\\
y=a_2+su_2+tv_2,\\
z=a_3+su_3+tv_3.
\end{cases}
\]

Como $u\times v$ es perpendicular a ambas direcciones, puede utilizarse como vector normal del plano.

\begin{definition}[Ecuación cartesiana de un plano]
Un plano de $\R^3$ puede escribirse
\[
\pi:\quad Ax+By+Cz+D=0,
\qquad (A,B,C)\neq(0,0,0).
\]
El vector
\[
n=(A,B,C)
\]
es un vector normal al plano.
\end{definition}

Si conocemos un punto $A=(a_1,a_2,a_3)$ y un vector normal $n=(n_1,n_2,n_3)$, la ecuación puede obtenerse directamente mediante
\[
n\cdot(X-A)=0,
\]
es decir,
\[
n_1(x-a_1)+n_2(y-a_2)+n_3(z-a_3)=0.
\]

Tres puntos no alineados $A,B,C$ determinan un plano. En este caso podemos tomar
\[
u=\overrightarrow{AB},
\qquad
v=\overrightarrow{AC},
\]
y calcular
\[
n=u\times v.
\]

\begin{solvedexercise}
Una superficie plana de captación está determinada, en un sistema de coordenadas local, por los puntos
\[
A=(0,0,2),\qquad B=(4,0,3),\qquad C=(0,3,4).
\]
Obtener una ecuación cartesiana del plano y un vector normal unitario.
\tcblower
\textbf{Resolución.}\par\smallskip
Tomamos
\[
u=B-A=(4,0,1),
\qquad
v=C-A=(0,3,2).
\]
Un vector normal es
\[
n=u\times v
=
\begin{vmatrix}
\mathbf i&\mathbf j&\mathbf k\\
4&0&1\\
0&3&2
\end{vmatrix}
=(-3,-8,12).
\]
Por tanto,
\[
-3(x-0)-8(y-0)+12(z-2)=0,
\]
y una ecuación cartesiana es
\[
\boxed{-3x-8y+12z-24=0.}
\]
Como
\[
\|n\|=\sqrt{9+64+144}=\sqrt{217},
\]
un vector normal unitario es
\[
\boxed{\widehat n=\frac1{\sqrt{217}}(-3,-8,12).}
\]
La orientación de una superficie puede describirse mediante su vector normal; por eso este tipo de cálculo aparece al estudiar superficies de montaje, paneles, placas y caras de elementos estructurales.
\end{solvedexercise}

\subsection{Posiciones relativas e intersecciones}

Las posiciones relativas se estudian combinando información sobre direcciones con la compatibilidad de sistemas de ecuaciones. Este punto conecta directamente la geometría con los capítulos de matrices y sistemas lineales.

\paragraph{Dos rectas en el plano.}
Sean
\[
r:A_1x+B_1y+C_1=0,
\qquad
s:A_2x+B_2y+C_2=0.
\]
El sistema formado por ambas ecuaciones puede ser:
\begin{itemize}
\item compatible determinado: las rectas son secantes;
\item compatible indeterminado: las rectas son coincidentes;
\item incompatible: las rectas son paralelas distintas.
\end{itemize}

Equivalentemente, si $v_r$ y $v_s$ son vectores directores, entonces las rectas son paralelas cuando $v_r$ y $v_s$ son proporcionales.

\paragraph{Recta y plano en el espacio.}
Sean
\[
r:X=A+tv,
\qquad
\pi:n\cdot(X-P)=0.
\]
La cantidad $n\cdot v$ determina la situación:
\begin{itemize}
\item si $n\cdot v\neq0$, la recta corta al plano en un único punto;
\item si $n\cdot v=0$ y $A\in\pi$, la recta está contenida en el plano;
\item si $n\cdot v=0$ y $A\notin\pi$, la recta es paralela al plano y no lo corta.
\end{itemize}

\begin{example}
Sea
\[
r:(x,y,z)=(1,0,2)+t(2,1,-1)
\]
y
\[
\pi: x-2y+z-5=0.
\]
El vector normal del plano es
\[
n=(1,-2,1).
\]
Como
\[
n\cdot v=(1,-2,1)\cdot(2,1,-1)=2-2-1=-1\neq0,
\]
la recta es secante al plano. Sustituyendo las ecuaciones paramétricas
\[
x=1+2t,\quad y=t,\quad z=2-t
\]
en la ecuación del plano,
\[
(1+2t)-2t+(2-t)-5=0,
\]
de donde
\[
t=-2.
\]
El punto de intersección es
\[
\boxed{(-3,-2,4).}
\]
\end{example}

\paragraph{Dos planos.}
Sean
\[
\pi_1:n_1\cdot X+d_1=0,
\qquad
\pi_2:n_2\cdot X+d_2=0.
\]
Si $n_1$ y $n_2$ no son proporcionales, los planos se cortan en una recta cuyo vector director puede tomarse como
\[
v=n_1\times n_2.
\]
Si los vectores normales son proporcionales, los planos son paralelos o coincidentes; para distinguir ambos casos basta comprobar si un punto de uno de ellos pertenece al otro.

\paragraph{Dos rectas en el espacio.}
Sean
\[
r:A+tu,
\qquad
s:B+\lambda v.
\]
Si $u$ y $v$ son proporcionales, las rectas son paralelas o coincidentes. Si no son proporcionales, resolvemos
\[
A+tu=B+\lambda v.
\]
Si el sistema tiene solución, las rectas son secantes. Si no tiene solución, son \emph{rectas cruzadas}: no son paralelas, pero tampoco se cortan.

Una caracterización útil de la coplanaridad es
\[
(B-A)\cdot(u\times v)=0.
\]
Cuando $u$ y $v$ no son paralelos, las rectas son secantes si y sólo si son coplanarias; si el producto mixto anterior es no nulo, son cruzadas.

\begin{solvedexercise}
Estudiar la posición relativa de
\[
r:(x,y,z)=(1,0,1)+t(1,2,-1)
\]
y
\[
s:(x,y,z)=(0,1,2)+\lambda(2,1,1).
\]
\tcblower
\textbf{Resolución.}\par\smallskip
Los vectores directores
\[
u=(1,2,-1),\qquad v=(2,1,1)
\]
no son proporcionales. Buscamos una posible intersección resolviendo
\[
(1,0,1)+t(1,2,-1)
=(0,1,2)+\lambda(2,1,1).
\]
Esto equivale a
\[
\begin{cases}
1+t=2\lambda,\\
2t=1+\lambda,\\
1-t=2+\lambda.
\end{cases}
\]
De la segunda ecuación,
\[
\lambda=2t-1.
\]
Sustituyendo en la primera,
\[
1+t=4t-2
\quad\Longrightarrow\quad
 t=1,
\]
y entonces $\lambda=1$. Sin embargo, la tercera ecuación daría
\[
0=3,
\]
lo cual es imposible. Por tanto, las rectas no se cortan.

Como tampoco son paralelas,
\[
\boxed{r\text{ y }s\text{ son rectas cruzadas}.}
\]
\end{solvedexercise}

\section{Problemas métricos}
\label{sec:problemas-metricos}

\subsection{Distancias}

La distancia entre objetos geométricos se define como la menor distancia entre sus puntos. Las fórmulas que utilizaremos proceden de proyecciones ortogonales y productos vectoriales.

\paragraph{Distancia de un punto a una recta del plano.}
Sea
\[
r:Ax+By+C=0
\]
y sea $P=(p_1,p_2)$. Entonces
\[
\boxed{
 d(P,r)=
 \frac{|Ap_1+Bp_2+C|}{\sqrt{A^2+B^2}}.}
\]

Esta expresión puede interpretarse proyectando el vector que une $P$ con un punto cualquiera de la recta sobre el vector normal $n=(A,B)$.

\begin{example}
Calcular la distancia del punto $P=(3,-1)$ a la recta
\[
r:2x-y-4=0.
\]
Aplicando la fórmula,
\[
d(P,r)
=\frac{|2\cdot3-(-1)-4|}{\sqrt{2^2+(-1)^2}}
=\frac{3}{\sqrt5}.
\]
Por tanto,
\[
\boxed{d(P,r)=\frac{3}{\sqrt5}.}
\]
\end{example}

\paragraph{Distancia de un punto a una recta del espacio.}
Sea
\[
r:X=A+tv
\]
y sea $P\in\R^3$. El área del paralelogramo generado por $P-A$ y $v$ es
\[
\|(P-A)\times v\|.
\]
También es igual a base por altura,
\[
\|v\|\,d(P,r).
\]
Por tanto,
\[
\boxed{
 d(P,r)=\frac{\|(P-A)\times v\|}{\|v\|}.}
\]

\begin{solvedexercise}
Calcular la distancia del punto
\[
P=(1,-2,3)
\]
a la recta
\[
r:(x,y,z)=(3,-1,0)+t(-1,1,1).
\]
\tcblower
\textbf{Resolución.}\par\smallskip
Tomamos
\[
A=(3,-1,0),
\qquad
v=(-1,1,1).
\]
Entonces
\[
P-A=(-2,-1,3)
\]
y
\[
(P-A)\times v
=
\begin{vmatrix}
\mathbf i&\mathbf j&\mathbf k\\
-2&-1&3\\
-1&1&1
\end{vmatrix}
=(-4,-1,-3).
\]
Por tanto,
\[
d(P,r)
=\frac{\sqrt{16+1+9}}{\sqrt{1+1+1}}
=\boxed{\sqrt{\frac{26}{3}}}.
\]
\end{solvedexercise}

\paragraph{Distancia de un punto a un plano.}
Sea
\[
\pi:Ax+By+Cz+D=0
\]
y $P=(p_1,p_2,p_3)$. Entonces
\[
\boxed{
 d(P,\pi)
 =\frac{|Ap_1+Bp_2+Cp_3+D|}{\sqrt{A^2+B^2+C^2}}.}
\]

De nuevo, el numerador mide ---salvo el factor $\|n\|$--- la componente del desplazamiento en la dirección normal al plano.

\begin{example}
Una carcasa metálica ocupa el plano
\[
\pi:2x-y+2z-6=0.
\]
Un sensor está situado en
\[
P=(3,1,4).
\]
La separación perpendicular entre el sensor y la superficie es
\[
d(P,\pi)
=\frac{|2\cdot3-1+2\cdot4-6|}{\sqrt{2^2+(-1)^2+2^2}}
=\frac{7}{3}.
\]
La distancia mínima es, por tanto,
\[
\boxed{\frac73}
\]
en las unidades del sistema de coordenadas utilizado.
\end{example}

\paragraph{Distancias entre objetos paralelos.}
Si dos rectas del plano son paralelas,
\[
r:Ax+By+C_1=0,
\qquad
s:Ax+By+C_2=0,
\]
entonces
\[
\boxed{
 d(r,s)=\frac{|C_2-C_1|}{\sqrt{A^2+B^2}}.}
\]
Análogamente, para dos planos paralelos
\[
\pi_1:Ax+By+Cz+D_1=0,
\qquad
\pi_2:Ax+By+Cz+D_2=0,
\]
se tiene
\[
\boxed{
 d(\pi_1,\pi_2)
 =\frac{|D_2-D_1|}{\sqrt{A^2+B^2+C^2}}.}
\]

\paragraph{Distancia entre dos rectas cruzadas.}
Sean
\[
r:A+tu,
\qquad
s:B+\lambda v,
\]
con $u\times v\neq0$. Si las rectas son cruzadas, el vector $u\times v$ es perpendicular a ambas direcciones. La distancia es la proyección de $B-A$ sobre esa dirección común:
\[
\boxed{
 d(r,s)
 =\frac{|(B-A)\cdot(u\times v)|}{\|u\times v\|}.}
\]
Si las rectas son secantes, el numerador vale cero y la fórmula proporciona, correctamente, distancia nula.

\subsection{Ángulos entre rectas y planos}

La medida de ángulos se reduce al producto escalar entre vectores directores o vectores normales.

\paragraph{Ángulo entre dos rectas.}
Si $u$ y $v$ son vectores directores de dos rectas, el ángulo geométrico entre ellas se toma habitualmente como el menor de los dos posibles, por lo que
\[
\boxed{
\cos\theta
=\frac{|u\cdot v|}{\|u\|\,\|v\|},
\qquad
0\leq\theta\leq\frac\pi2.}
\]

\paragraph{Ángulo entre dos planos.}
Si $n_1$ y $n_2$ son vectores normales de los planos, el ángulo entre los planos es el ángulo agudo entre sus normales:
\[
\boxed{
\cos\theta
=\frac{|n_1\cdot n_2|}{\|n_1\|\,\|n_2\|}.}
\]

\paragraph{Ángulo entre una recta y un plano.}
Sea $v$ un vector director de la recta y $n$ un vector normal al plano. Si $\alpha$ es el ángulo entre la recta y el plano, entonces
\[
\alpha+\theta=\frac\pi2,
\]
donde $\theta$ es el ángulo entre $v$ y $n$. Por tanto,
\[
\boxed{
\sin\alpha
=\frac{|v\cdot n|}{\|v\|\,\|n\|}.}
\]

\begin{solvedexercise}
Una conducción rectilínea tiene dirección
\[
v=(2,-1,2)
\]
y atraviesa una pared cuyo plano es
\[
\pi:x+2y+2z-8=0.
\]
Calcular el ángulo que forma la conducción con la pared.
\tcblower
\textbf{Resolución.}\par\smallskip
Un vector normal al plano es
\[
n=(1,2,2).
\]
Entonces
\[
v\cdot n=2-2+4=4,
\]
\[
\|v\|=3,
\qquad
\|n\|=3.
\]
Por tanto,
\[
\sin\alpha=\frac49,
\]
y
\[
\boxed{\alpha=\arcsin\left(\frac49\right).}
\]
Esta es la inclinación de la conducción respecto de la superficie, no respecto de la normal a la superficie.
\end{solvedexercise}

Las condiciones de paralelismo y perpendicularidad se obtienen como casos particulares. Por ejemplo:
\begin{itemize}
\item dos rectas son perpendiculares si sus vectores directores tienen producto escalar cero;
\item dos planos son perpendiculares si sus normales tienen producto escalar cero;
\item una recta es perpendicular a un plano si su vector director es paralelo al vector normal del plano;
\item una recta es paralela a un plano si su vector director es ortogonal al normal del plano.
\end{itemize}

\subsection{Proyecciones y puntos de mínima distancia}

Las fórmulas de distancia proporcionan un número, pero en muchos problemas interesa además localizar el punto que realiza esa distancia mínima.

\paragraph{Proyección de un punto sobre una recta.}
Sea
\[
r:X=A+tv
\]
y sea $P$ un punto. El pie de la perpendicular $H$ es un punto de la recta,
\[
H=A+t_0v,
\]
y debe satisfacer
\[
(P-H)\cdot v=0.
\]
Sustituyendo,
\[
(P-A-t_0v)\cdot v=0,
\]
de donde
\[
\boxed{
 t_0=\frac{(P-A)\cdot v}{v\cdot v}.}
\]
Así,
\[
\boxed{
 H=A+\frac{(P-A)\cdot v}{v\cdot v}\,v.}
\]

\paragraph{Proyección de un punto sobre un plano.}
Sea
\[
\pi:n\cdot(X-A)=0
\]
y $P$ un punto. El punto proyectado $H$ debe encontrarse sobre la recta perpendicular al plano que pasa por $P$:
\[
H=P-tn.
\]
Imponiendo $H\in\pi$,
\[
n\cdot(P-tn-A)=0,
\]
por lo que
\[
t=\frac{n\cdot(P-A)}{n\cdot n}.
\]
En consecuencia,
\[
\boxed{
 H=P-\frac{n\cdot(P-A)}{n\cdot n}\,n.}
\]

\begin{solvedexercise}
Una herramienta se encuentra en
\[
P=(4,1,5)
\]
y debe aproximarse perpendicularmente a la superficie
\[
\pi:2x-y+2z-7=0.
\]
Determinar el punto $H$ de la superficie al que llega siguiendo la trayectoria más corta y la longitud de dicha trayectoria.
\tcblower
\textbf{Resolución.}\par\smallskip
Tomamos
\[
n=(2,-1,2).
\]
Podemos utilizar cualquier punto del plano; por ejemplo, $A=(0,-7,0)$ satisface
\[
2\cdot0-(-7)+2\cdot0-7=0.
\]
Entonces
\[
P-A=(4,8,5),
\]
y
\[
n\cdot(P-A)=8-8+10=10,
\qquad
n\cdot n=9.
\]
Por tanto,
\[
H=P-\frac{10}{9}(2,-1,2)
=\left(\frac{16}{9},\frac{19}{9},\frac{25}{9}\right).
\]
La longitud de la trayectoria es
\[
\|P-H\|
=\frac{10}{9}\|n\|
=\frac{10}{3}.
\]
Así,
\[
\boxed{
H=\left(\frac{16}{9},\frac{19}{9},\frac{25}{9}\right),
\qquad
d(P,\pi)=\frac{10}{3}.}
\]
Como comprobación, al sustituir $H$ en la ecuación del plano se obtiene cero.
\end{solvedexercise}

Los problemas métricos suelen requerir combinar varias de las ideas del capítulo. Antes de utilizar una fórmula conviene identificar con precisión qué objetos intervienen y qué información los determina: puntos, vectores directores y normales. Después se traduce el problema a productos escalares, productos vectoriales o sistemas lineales.

\begin{example}[Distancia entre dos ejes no coplanarios]
Dos ejes de una instalación se modelizan mediante
\[
r:(x,y,z)=(0,0,0)+t(1,0,1),
\]
\[
s:(x,y,z)=(1,2,0)+\lambda(0,1,1).
\]
Como
\[
u=(1,0,1),\qquad v=(0,1,1),
\]
no son paralelos. Además,
\[
u\times v=(-1,-1,1).
\]
Con $A=(0,0,0)$ y $B=(1,2,0)$,
\[
(B-A)\cdot(u\times v)
=(1,2,0)\cdot(-1,-1,1)=-3\neq0,
\]
por lo que las rectas son cruzadas. Su distancia es
\[
d(r,s)
=\frac{3}{\sqrt3}
=\boxed{\sqrt3}.
\]
El mismo producto mixto que detecta que las rectas no son coplanarias proporciona, tras normalizar, la separación mínima entre ellas.
\end{example}

\section*{Ejercicios propuestos}
\addcontentsline{toc}{section}{Ejercicios propuestos}

Los ejercicios siguientes combinan cálculo vectorial, obtención de ecuaciones, estudio de posiciones relativas y problemas métricos. Algunos proceden del tipo de ejercicios utilizados en cursos anteriores y otros amplían la colección para trabajar de forma integrada los contenidos del capítulo. Los marcados con \M son especialmente adecuados para ser resueltos o comprobados también con \textit{Wolfram Mathematica}.

\begin{enumerate}

\item Dados los puntos
\[
A=(1,-4),\qquad B=(5,7),
\]
calcular $\overrightarrow{AB}$, la distancia $d(A,B)$ y un vector unitario con la dirección de $A$ hacia $B$.

\item Dados
\[
P=(1,2,3),\qquad Q=(5,-3,1),
\]
calcular la distancia entre ambos puntos. Determinar además el punto medio del segmento $PQ$.

\item Para
\[
u=(2,-1,3),\qquad v=(1,4,-2),
\]
calcular:
\begin{enumerate}[label=\alph*)]
\item $u\cdot v$;
\item las normas de $u$ y $v$;
\item el ángulo que forman;
\item la proyección de $u$ sobre $v$;
\item una descomposición $u=u_{\parallel}+u_{\perp}$ respecto de la dirección de $v$.
\end{enumerate}

\item Determinar el valor de $a\in\R$ para que los vectores
\[
u=(1,a,2),\qquad v=(2,-1,1)
\]
sean ortogonales. Para ese valor, calcular el área del paralelogramo que generan.

\item Dados
\[
u=(1,2,-1),\qquad v=(2,0,3),
\]
calcular $u\times v$ y comprobar que el resultado es ortogonal a ambos vectores.

\item Calcular el área del triángulo de vértices
\[
A=(1,0,2),\qquad B=(3,-1,1),\qquad C=(0,2,4).
\]

\item Estudiar si los vectores
\[
u=(1,2,0),\qquad v=(0,1,1),\qquad w=(2,5,1)
\]
son coplanarios. En caso contrario, calcular el volumen del paralelepípedo que generan.

\item Un brazo mecánico aplica una fuerza
\[
F=(80,30,-40)\ \mathrm{N}
\]
en un punto cuyo vector de posición respecto del eje de giro es
\[
r=(0.25,0.10,0)\ \mathrm{m}.
\]
Calcular el momento $M=r\times F$ y su módulo.

\item Una corriente de fluido tiene velocidad
\[
v=(3,4,0)\ \mathrm{m/s}
\]
y una sección de control posee vector normal unitario
\[
n=\frac{1}{3}(2,-1,2).
\]
Calcular la componente de la velocidad normal a la sección y la componente tangencial.

\item Hallar las ecuaciones vectorial y paramétricas de la recta del plano que pasa por
\[
P=(2,-1),\qquad Q=(-1,5),
\]
y obtener después una ecuación cartesiana.

\item Hallar la ecuación de la recta que pasa por $P=(3,1)$ y es perpendicular a
\[
r:2x-3y+4=0.
\]

\item Estudiar la posición relativa de las rectas
\[
r:2x-y+1=0,
\qquad
s:4x-2y-5=0.
\]
Calcular la distancia entre ellas cuando proceda.

\item Hallar las ecuaciones vectorial, paramétricas y continuas de la recta del espacio que pasa por
\[
P=(1,0,-1),\qquad Q=(2,1,-3).
\]

\item Hallar la recta que pasa por $A=(1,-3,0)$ y es paralela al vector $u\times v$, siendo
\[
u=(-1,-1,2),\qquad v=(6,0,0).
\]

\item La trayectoria ideal del extremo de un robot pasa por el punto
\[
A=(0.5,1,0)
\]
y tiene dirección
\[
v=(2,-1,3).
\]
Determinar el punto alcanzado para $t=2$ y decidir si
\[
P=(4.5,-1,6)
\]
pertenece a la trayectoria.

\item Obtener una ecuación cartesiana del plano que pasa por
\[
P=(0,1,-1),\qquad Q=(2,3,-5),\qquad R=(1,4,3).
\]

\item Hallar la ecuación del plano cuyo vector normal es
\[
n=(-1,2,-3)
\]
y que pasa por
\[
A=(2,0,4).
\]

\item Hallar las ecuaciones paramétricas y cartesianas del plano determinado por
\[
A=(1,-3,2),\qquad u=(2,1,0),\qquad v=(-1,0,3).
\]

\item Determinar el valor de $a$ para que el punto
\[
P=(1,a,2)
\]
pertenezca al plano
\[
\pi:2x-y+3z-5=0.
\]
Para ese valor, obtener una recta perpendicular a $\pi$ que pase por $P$.

\item Una placa de intercambio térmico está contenida en el plano que pasa por los puntos
\[
A=(0,0,0),\qquad B=(2,1,0),\qquad C=(0,2,1).
\]
Obtener un vector normal unitario a la placa y su ecuación cartesiana.

\item Estudiar la posición relativa de la recta
\[
r:(x,y,z)=(1,2,-1)+t(1,-1,2)
\]
y el plano
\[
\pi:2x+y-z-4=0.
\]
Si se cortan, calcular el punto de intersección.

\item Estudiar la posición relativa de
\[
r:(x,y,z)=(1,0,0)+t(1,2,1)
\]
y
\[
s:(x,y,z)=(0,1,1)+\lambda(2,4,2).
\]

\item Estudiar la posición relativa de
\[
r:(x,y,z)=(1,0,1)+t(1,1,0)
\]
y
\[
s:(x,y,z)=(0,1,2)+\lambda(0,1,1).
\]
Si son secantes, calcular el punto de corte; si son cruzadas, justificarlo mediante un producto mixto.

\item Estudiar la posición relativa de los planos
\[
\pi_1:x+y-z+1=0,
\qquad
\pi_2:2x-y+z-3=0.
\]
Si se cortan, obtener una ecuación paramétrica de la recta de intersección.

\item Estudiar, según el valor del parámetro $a$, la posición relativa de los planos
\[
\pi_1:x+2y-z=1,
\qquad
\pi_2:2x+4y-2z=a.
\]

\item Calcular la distancia del punto
\[
P=(2,-3)
\]
a la recta
\[
r:3x+4y-7=0.
\]
Obtener además el pie de la perpendicular desde $P$ a $r$.

\item Calcular la distancia del punto
\[
P=(1,-2,3)
\]
a la recta
\[
r:(x,y,z)=(3,-1,0)+t(-1,1,1).
\]
Obtener el punto de la recta más próximo a $P$.

\item Calcular la distancia del punto
\[
P=(1,2,3)
\]
al plano
\[
\pi:2x+3y-z=-7.
\]
Obtener también la proyección ortogonal de $P$ sobre $\pi$.

\item Calcular la distancia entre los planos paralelos
\[
\pi_1:2x-y+2z-3=0,
\qquad
\pi_2:4x-2y+4z+5=0.
\]
Antes de aplicar una fórmula, escribir ambos planos con el mismo vector normal.

\item Calcular la distancia entre las rectas
\[
r:(x,y,z)=(0,0,0)+t(1,1,0),
\]
\[
s:(x,y,z)=(1,0,2)+\lambda(1,-1,0).
\]
Estudiar primero su posición relativa.

\item Hallar el ángulo agudo entre las rectas cuyas direcciones son
\[
u=(1,2,2),\qquad v=(2,-1,2).
\]

\item Hallar el ángulo entre los planos
\[
\pi_1:x+2y-2z+1=0,
\qquad
\pi_2:2x-y+2z-3=0.
\]

\item Hallar el ángulo que forma la recta
\[
r:(x,y,z)=(1,0,2)+t(2,1,-2)
\]
con el plano
\[
\pi:x-y+2z=0.
\]

\item Una conducción debe atravesar perpendicularmente una pared
\[
\pi:2x-y+2z-9=0.
\]
Si el punto de partida es
\[
P=(1,2,0),
\]
obtener la recta de perforación, el punto donde alcanza la pared y la distancia recorrida hasta ella.

\item Una cámara de inspección está situada en
\[
P=(5,2,3)
\]
y debe orientarse hacia el punto más próximo del eje de una tubería
\[
r:(x,y,z)=(1,0,1)+t(2,1,0).
\]
Determinar dicho punto, la distancia y un vector unitario que indique la dirección desde la cámara hasta la tubería.

\item En un sistema eléctrico tridimensional, dos tramos rectilíneos de cableado idealizados tienen ejes
\[
r:(x,y,z)=(0,0,1)+t(1,2,0),
\]
\[
s:(x,y,z)=(2,0,0)+\lambda(0,1,2).
\]
Estudiar si los ejes se cortan. Si son cruzados, calcular su separación mínima.

\item En una instalación energética, la superficie de un panel tiene normal
\[
n=(1,-2,2)
\]
y la dirección de incidencia de un haz se modeliza mediante
\[
v=(-2,1,-2).
\]
Calcular el ángulo entre el haz y la superficie. Determinar también el ángulo entre el haz y la normal.

\item Una tubería de proceso tiene eje
\[
r:(x,y,z)=(2,1,0)+t(1,-2,2),
\]
y debe atravesar el plano
\[
\pi:x+y+z-6=0.
\]
Determinar el punto de paso y el ángulo de incidencia. Interpretar qué ocurriría si el vector director se modificase a $v=(1,-2,1)$.

\item \M Comprobar con Mathematica los ejercicios 3--9 utilizando únicamente productos escalares, \texttt{Norm} y \texttt{Cross} cuando corresponda. Para el ángulo entre dos vectores puede evaluarse
\[
\operatorname{ArcCos}\left(\frac{u\cdot v}{\|u\|\|v\|}\right).
\]
Comparar siempre el resultado con el cálculo manual.

\item \M En los ejercicios 10--25, utilizar \texttt{Solve} para comprobar pertenencias e intersecciones. Plantear primero manualmente las ecuaciones y utilizar el ordenador únicamente para resolver o verificar el sistema resultante.

\item \M En los ejercicios 26--38, comprobar las distancias y proyecciones con combinaciones de \texttt{Norm}, producto escalar y \texttt{Cross}. No introducir comandos específicos de geometría: el objetivo es reutilizar las mismas operaciones vectoriales estudiadas en la teoría.

\item (Examen) Dados
\[
r:(x,y,z)=(1,-1,2)+t(1,2,-1)
\]
y
\[
\pi:2x-y+z-4=0,
\]
se pide:
\begin{enumerate}[label=\alph*)]
\item estudiar la posición relativa de $r$ y $\pi$;
\item calcular, si existe, su punto de intersección;
\item hallar el ángulo entre la recta y el plano;
\item calcular la distancia del punto $P=(3,0,1)$ al plano y su proyección ortogonal sobre él.
\end{enumerate}

\item (Examen) Sean las rectas
\[
r:(x,y,z)=(1,0,0)+t(1,1,2),
\]
\[
s:(x,y,z)=(0,2,1)+\lambda(2,-1,1).
\]
\begin{enumerate}[label=\alph*)]
\item estudiar su posición relativa;
\item calcular la distancia entre ambas;
\item obtener un vector perpendicular a las dos direcciones;
\item si son cruzadas, explicar geométricamente por qué el vector que realiza la distancia mínima debe ser paralelo a ese vector perpendicular común.
\end{enumerate}

\item (Examen) Una superficie plana queda determinada por
\[
A=(1,0,1),\qquad B=(3,1,0),\qquad C=(0,2,2).
\]
Una trayectoria rectilínea parte de
\[
P=(2,-1,4)
\]
con dirección $v=(1,1,-2)$. Se pide:
\begin{enumerate}[label=\alph*)]
\item obtener la ecuación cartesiana del plano $ABC$;
\item estudiar la posición relativa de la trayectoria y el plano;
\item calcular, si existe, el punto de intersección;
\item calcular el ángulo entre la trayectoria y el plano;
\item obtener la distancia de $P$ al plano y el punto de la superficie más próximo a $P$.
\end{enumerate}

\end{enumerate}
